\documentclass[reprint]{revtex4-2}
\usepackage{amsmath,amssymb,graphicx,hyperref,booktabs}
\begin{document}

\title{Information Capacity and the Quantum-Classical Separation: \\
Why the Universe Looks Different at Different Scales}

\author{Huang Zhongchang}
\affiliation{Independent Researcher, Nanjing, China}

\begin{abstract}
We show that the separation between quantum and classical behavior follows from three information-theoretic axioms: distinguishable states exist (A1), capacity is bounded (A2), and dynamics is local (A3). From these axioms we prove a chain of theorems---bijectivity of microdynamics, coarse-grained entropy identity, Kac recurrence scaling, and an information reflux bound---that together establish a continuum $q$-field whose static limit satisfies $\nabla^2(1/q)=0$, structurally identical to the Laplace equation for the Newtonian gravitational potential. The $q\to 0$ regime corresponds to a self-organized critical state: cells are nearly saturated, any perturbation triggers cascaded overflow, and the static geometry of this critical state is what we measure as gravity. The $q\approx 1/2$ regime is the statistical preference of the mixing entropy, explaining why quantum coherence persists. The framework makes one clean experimental prediction: an information reflux upper bound $P_{\rm reflux}\le q_S/q_E$, testable in circuit QED at zero temperature where standard Lindblad theory predicts strictly zero reflux. We do not claim to derive the universe---we claim to show that \emph{given} the three axioms, the separation of scales, the arrow of time, and the structure of the gravitational field equation are not independent postulates but consequences of a single underlying principle: information flow under finite capacity.
\end{abstract}

% ============================================================
\section{Introduction}
% ============================================================

Why does the universe appear quantum at microscopic scales and classical at macroscopic scales? Standard decoherence theory~\cite{Zurek2003,Schlosshauer2007} answers \emph{how} this transition occurs: environmental monitoring selects pointer states, and the system's reduced density matrix becomes diagonal in that basis. Quantum Darwinism~\cite{Zurek2009} further explains how objectivity emerges from redundant environmental records.

But standard theory does not answer a deeper question: \emph{why} does the separation happen in this direction? Why does the classical world grow at the expense of the quantum world, rather than the reverse? Why is the separation irreversible? And why does the gravitational field equation share the mathematical structure of an information diffusion equation in its static limit?

This paper proposes that the answer lies not in dynamics but in the logical structure of information itself. We begin with three axioms that define what it means for a system to ``exist'' and ``interact'':

\begin{itemize}
\item \textbf{A1 (Distinguishability):} States that differ can be distinguished. Without this, ``existence'' has no operational meaning---nothing can be identified as being one thing rather than another.
\item \textbf{A2 (Bounded capacity):} The total number of distinguishable states is finite. Without this, ``interaction'' cannot be defined---there is no boundary between system and environment, and ``information transfer'' becomes an untraceable claim.
\item \textbf{A3 (Locality):} The dynamics is $c$-local with respect to some cell decomposition: each output cell depends on at most $c$ input cells, with $c\ll N$. This encodes the physical fact that interactions are local---a consequence of finite propagation speed, though the framework does not presuppose spacetime.
\end{itemize}

From these axioms we prove that:
\begin{enumerate}
\item The microdynamics is a bijection on the state space (Theorem~1), implying information is structurally conserved but no specific macroscopic quantity need be.
\item Coarse-grained Shannon entropy change equals the change in intra-fiber KL divergence (Theorem~2), clarifying the statistical origin of the second law.
\item Kac recurrence time scales as $2^{N_S}$, independent of environment size $N_E$ (Theorem~3), distinguishing relaxation (controlled by $N_E$) from recurrence (controlled by $N_S$).
\item Information reflux from environment to system satisfies $P_{\rm reflux}\le q_S/q_E$ (Theorem~S1), where $q$ is the vacant capacity fraction.
\end{enumerate}

Taking the continuum limit yields a $q$-field obeying $\partial_t q = D_0\nabla^2(1/q)$ (Theorems~4--7). The $q\to 0$ limit reveals a self-organized critical state---cells are nearly saturated, any perturbation triggers cascaded overflow, and the static geometry satisfies $\nabla^2(1/q)=0$, the Laplace equation for the gravitational potential (Theorems~11--12).

The framework makes one clean experimental prediction: at zero temperature, standard Lindblad theory predicts zero information reflux from environment to system, while our reflux bound predicts a finite upper limit $P_{\rm reflux}\le q_S/q_E$. This can be tested in circuit QED with existing technology.

We do not claim to \emph{derive} the universe from first principles. We claim something more modest but arguably deeper: given three axioms about information, the separation of quantum and classical scales, the arrow of time, and the structure of the gravitational field equation emerge as consequences of a single principle---information flow under finite capacity.

% ============================================================
\section{Axiomatic Framework}
% ============================================================

\subsection{The Cell Model}

Consider a system of $N$ binary cells. The global state is $\mathbf{s}=(s_1,\ldots,s_N)\in\{0,1\}^N$. Each cell $i$ has state $s_i=0$ (vacant, capable of receiving information) or $s_i=1$ (occupied, at capacity). The total state space is $X=\{0,1\}^N$, with $|X|=2^N$~bits.

Dynamics is a map $f:X\to X$ taking the system from one state to the next. Axiom~A1 requires $f$ to be injective: distinguishable states remain distinguishable after evolution. Axiom~A2 fixes $N$ finite. Together they imply:

\begin{theorem}[Bijection]
\label{thm:bijection}
For any finite set $X$ and injective map $f:X\to X$, $f$ is bijective.
\end{theorem}
\begin{proof}
An injective map from a finite set to itself is surjective by the pigeonhole principle. $\square$
\end{proof}

This is the structural conservation law of the framework: the total \emph{number of distinguishable states} is conserved, because $f$ permutes $X$. However---and this is crucial---no specific macroscopic quantity need be conserved:

\begin{theorem}[Non-conservation of occupancy]
\label{thm:nonconservation}
Axioms~A1--A2 do \emph{not} imply conservation of the total occupation number $\Sigma s_i$.
\end{theorem}
\begin{proof}
Counterexample: $N=2$, $f(00)=01$, $f(01)=10$, $f(10)=11$, $f(11)=00$. This is a bijection on $\{0,1\}^2$. The Hamming weight $|s|=\Sigma s_i$ evolves as $0\to 1\to 1\to 2\to 0$, which is not constant. $\square$
\end{proof}

This theorem corrects an imprecise claim in earlier versions of this framework: ``information conservation'' means structural conservation (bijectivity), not conservation of any specific macroscopic quantity.

\subsection{The $q$ Parameter}

Define the macroscopic vacant fraction:
\begin{equation}
q(\mathbf{s}) = \frac{N-\sum_i s_i}{N}\in[0,1].
\end{equation}
$q=1$: all cells vacant (maximal capacity, fully quantum). $q=0$: all cells occupied (zero capacity, classical). Intermediate $q$ corresponds to partial occupation.

Axiom~A3 (locality) is needed to make the cell decomposition well-defined. Without it, a random bijection $f$ admits no $c$-local cell decomposition with probability approaching~1 for $c=o(N)$; A3 selects the physically relevant class of dynamics.

\subsection{Relation to Standard Quantum Mechanics}

In Section~\ref{sec:interpretation} we establish the precise mapping: $q$ is the normalized $\ell_1$-norm of coherence in the pointer basis:
\begin{equation}
q(\rho) = \frac{\sum_{i\neq j}|\rho_{ij}|}{\sum_{i,j}|\rho_{ij}|} = \frac{\|\rho-\mathrm{diag}(\rho)\|_1}{\|\rho\|_1}.
\end{equation}
This mapping makes DGF compatible with standard quantum mechanics: it is not a competing theory but an information-theoretic \emph{reformulation} that reveals constraints (such as the reflux bound) not visible in the standard formalism.

% ============================================================
\section{Discrete Theorems}
% ============================================================

\subsection{Coarse-Grained Entropy Identity}

Define the coarse-graining map $\pi:X\to[0,1]$ projecting microstates onto their $q$-value. The fiber $\pi^{-1}(q)$ contains $\binom{N}{N(1-q)}$ microstates. The macroscopic distribution $P(q)$ evolves under the induced transition matrix $W(q'|q)$.

\begin{theorem}[Entropy Identity]
\label{thm:entropy}
Let $\rho_q(\mathbf{s})$ be the distribution within fiber $q$, and $D[\rho\|\mu]$ the KL divergence from the uniform distribution on the fiber. Then:
\begin{equation}
S[P'] - S[P] = \mathbb{E}_{q'}[D[\rho'_{q'}\|\mu_{q'}]] - \mathbb{E}_q[D[\rho_q\|\mu_q]].
\end{equation}
\end{theorem}
\begin{proof}
By direct computation using $S[P] = -\sum_q P(q)\log P(q)$ and the definition of $W(q'|q)$ induced by the bijective $f$. The identity holds exactly for any bijection, with no thermodynamic limit required.
\end{proof}

This identity separates entropy change into two contributions: the coarse-grained part (which can increase or decrease) and the intra-fiber part (which measures deviation from uniformity). When the intra-fiber distribution is initially uniform ($D[\rho]=0$), entropy can only increase---but this is a \emph{lower bound from the initial condition}, not a monotonic stepwise law. Numerical verification on $N=3$ confirms that entropy can decrease between consecutive steps while remaining above its initial value.

\subsection{Kac Recurrence}

\begin{theorem}[Recurrence time]
\label{thm:kac}
For a subsystem of $N_S$ cells within a total of $N=N_S+N_E$ cells, under random bijective dynamics, the expected Kac recurrence time for the subsystem state is $\langle T_{\rm rec}\rangle = 2^{N_S}$, independent of $N_E$.
\end{theorem}
\begin{proof}
The subsystem state space has $2^{N_S}$ elements. By Kac's lemma, for any ergodic dynamical system, the expected recurrence time equals the inverse of the stationary measure: $\langle T_{\rm rec}\rangle = 1/\pi(s_S) = 2^{N_S}$. The environment size $N_E$ does not appear because the stationary distribution is uniform over all $2^{N_S}$ subsystem states regardless of $N_E$.
\end{proof}

This theorem cleanly separates two timescales: the \emph{relaxation} time (how fast the subsystem approaches equilibrium, controlled by $N_E$) and the \emph{recurrence} time (how long until an exact state repeats, controlled by $N_S$). The standard intuition---``large environment $\Rightarrow$ irreversible''---confuses fast relaxation with fundamental irreversibility. Theorem~3 shows that even an infinite environment cannot eliminate recurrence; it can only delay relaxation.

\subsection{Information Reflux Bound}

\begin{theorem}[Reflux bound, Case A: conserved total excitation]
\label{thm:reflux}
When the total excitation number $K=\sum_i s_i$ is conserved (a good approximation for many physical systems at low temperature), the probability of information reflux---a measured increase in subsystem excitation by $\delta$ or more within observation time $T$---satisfies:
\begin{equation}
P_{\rm reflux} \le T \cdot \left(\frac{q_S}{1-q_S}\right)^\delta \cdot \left(\frac{1-q_E}{q_E}\right)^\delta,
\end{equation}
where $q_S=1-k_{\rm min}/N_S$ and $q_E=1-(K-k_{\rm min})/N_E$ are the vacant fractions at the reflux starting point.
\end{theorem}
\begin{proof}
Under the uniform mixing property (UMP), macrostate transition probabilities are proportional to the number of accessible microstates. The transition $k\to k+\delta$ requires $\delta$ occupied cells in the environment to become vacant, each selecting from $N_E(1-q_E)$ occupied sites, while $\delta$ vacant cells in the system become occupied, each selecting from $N_S q_S$ vacant sites. The product of binomial factors gives the stated bound. The $\delta$-dependence captures the combinatorial suppression of large reflux events.
\end{proof}

For Case~B (total excitation not conserved), the reflux bound simplifies to $P_{\rm reflux}\le T\cdot\sum_{j}\binom{N_S}{j}/2^{N_S}$, independent of $q_E$---the environment provides space but not resistance. Physical systems typically fall between these extremes, with approximate conservation holding over experimentally relevant timescales.

% ============================================================
\section{Continuum $q$-Field}
% ============================================================

\subsection{From Discrete to Continuous}

Place $N$ cells on a $d$-dimensional lattice with spacing $a$. Define the local $q$-field as the coarse-grained vacant fraction over a mesoscopic volume. The discrete flux between neighboring cells $i$ and $j$ is:
\begin{equation}
J_{i\to j} = \frac{1}{q_i} - \frac{1}{q_j},
\end{equation}
which encodes the information-theoretic pressure: occupied cells (low $q$) exert higher ``pressure'' than vacant cells (high $q$). In the continuum limit $a\to 0$, $N\to\infty$ with fixed total volume:

\begin{theorem}[Continuum q-field equation]
\label{thm:continuum}
The $q$-field satisfies:
\begin{equation}
\boxed{\partial_t q = D_0\,\nabla^2\!\left(\frac{1}{q}\right)},
\end{equation}
with $D_0 = a^2/\tau_0 = a\cdot c$, where $\tau_0=a/c$ is the cell crossing time.
\end{theorem}

Equivalent forms---$\partial_t q = D_0\nabla\cdot(q^{-2}\nabla q)$, with effective diffusivity $D(q)=D_0/q^2$---clarify the nonlinearity. The equation conserves total $q$ and admits a Lyapunov functional $F_1[q]=\int q^{-1}d^dx$ (monotonically decreasing).

\subsection{The $q\to 0$ Singularity}

\begin{theorem}[Fast diffusion classification]
\label{thm:fastdiff}
Equation~(5) belongs to the fast diffusion class $m=-1$ (supercritical for $d\ge 2$). The effective diffusivity $D(q)=D_0/q^2$ diverges as $q\to 0$.
\end{theorem}

This divergence is the central physical feature of the $q$-field. When a region approaches $q\approx 0$, any local gradient produces divergent flux---the region becomes a powerful \emph{source} of information outflow.

\subsection{Physical Interpretation: Self-Organized Criticality}

The divergence $D\to\infty$ as $q\to 0$ is a \emph{continuum artifact}, not a physical prediction of infinite propagation speed. In the discrete model, information propagates at exactly $c$ per step. What the divergence correctly captures is the \emph{self-organized criticality} of the $q\to 0$ state:

\begin{itemize}
\item Cells in a $q\approx 0$ region are nearly saturated---each has minimal buffer capacity.
\item Any incoming perturbation finds almost no available capacity in the receiving cell.
\item The perturbation triggers immediate overflow to neighboring cells---which themselves have minimal buffer.
\item This cascaded overflow chain propagates with step time $\tau_0\to 0$ \emph{because there is no waiting time between steps}, not because individual steps exceed $c$.
\end{itemize}

The analogy is to a sandpile at the angle of repose~\cite{Bak1987}: adding one grain can trigger an avalanche that spans the entire pile. The avalanche does not propagate faster than the grain-to-grain interaction speed; it propagates with near-zero waiting time between interactions because every grain is already at threshold.

In the $q$-field, this criticality manifests as the divergent diffusivity. The physical cutoff is $q_{\rm min}\sim(a/L)^d$, below which the continuum description breaks down and the discrete cell dynamics must be used.

\subsection{Statistical Preference for $q\approx 1/2$}

\begin{theorem}[Entropic preference]
\label{thm:entropic}
The mixing entropy per cell $s(q)=-q\log q-(1-q)\log(1-q)$ has a unique global maximum at $q=1/2$. The free energy $f(q)=-T_0 s(q)$ has a unique global minimum at $q=1/2$.
\end{theorem}

This is the information-theoretic analog of the second law: the state space volume is maximal at $q=1/2$ (Stirling's approximation: $\binom{N}{N/2}\sim 2^N/\sqrt{\pi N/2}$), so statistical mechanics drives the system toward this point. The tension between the diffusive dynamics (which creates $q\approx 0$ cores and $q\approx 1/2$ backgrounds) and the entropic preference (which favors uniform $q\approx 1/2$) is the engine of the quantum-classical separation.

\subsection{Separation Mechanism}

The $q$-field dynamics together with the $q\to 0$ criticality and the entropic preference for $q\approx 1/2$ produce a two-domain asymptotic state:

\begin{theorem}[Two-domain asymptotics]
\label{thm:twodomain}
Under the $q$-field equation~(5) with initial condition containing spatial inhomogeneity, the asymptotic state consists of compact $q\approx 0$ cores (classical objects) embedded in a $q\approx 1/2$ sea (quantum background). The cores are stable: once a region reaches $q\approx 0$, the divergent diffusivity ensures it cannot return to higher $q$---any incoming information is immediately re-emitted.
\end{theorem}

The separation is not due to two worlds ``diverging.'' It is due to one regime ($q\approx 0$) becoming critically sensitive and the other ($q\approx 1/2$) remaining at statistical equilibrium. The boundary between them is sharp because the diffusivity contrast $D(0)/D(1/2)\to\infty$ prevents mixing.

% ============================================================
\section{Gravity from the Critical State}
% ============================================================

\subsection{Static Limit and the Laplace Equation}

\begin{theorem}[Static-gravity correspondence]
\label{thm:staticgravity}
In the static limit $\partial_t q=0$, the information potential $u\equiv 1/q$ satisfies:
\begin{equation}
\boxed{\nabla^2 u = 0},
\end{equation}
the Laplace equation---mathematically isomorphic to the vacuum Einstein equation in the Newtonian limit $\nabla^2\Phi=0$.
\end{theorem}
\begin{proof}
Set $\partial_t q=0$ in Eq.~(5): $D_0\nabla^2(1/q)=0$. Since $D_0>0$, $\nabla^2(1/q)=0$. Define $u\equiv 1/q$; then $\nabla^2 u=0$. $\square$
\end{proof}

With source terms (spatial regions where information is locked in occupied cells), the field obeys $\nabla^2 u = -\kappa\rho_m$, structurally identical to the Poisson equation $\nabla^2\Phi = 4\pi G\rho$. The coupling constant $\kappa$ is calibrated to $4\pi G/c^2$ by matching to Newtonian gravity, but the \emph{form} of the equation---Laplace/Poisson type---is derived, not assumed.

\subsection{Spherical Solution}

\begin{theorem}[Spherical weak-field solution]
\label{thm:spherical}
For a point mass $M$ in $d=3$ spatial dimensions, the static spherical solution with boundary condition $q(\infty)=1$ is:
\begin{equation}
\boxed{\frac{1}{q(r)} = 1 + \frac{GM}{rc^2}},
\end{equation}
in the weak-field regime $r\gg GM/c^2$. This yields the Newtonian correspondence $1/q = 1 - \Phi/c^2$.
\end{theorem}
\begin{proof}
In spherical symmetry, the Laplace equation is $(1/r^2)\partial_r(r^2\partial_r u)=0$. Integrating twice: $u(r) = -C_1/r + C_2$. Boundary conditions: $u(\infty)=1\Rightarrow C_2=1$; the mass $M$ is identified from the Gauss flux $\oint\nabla u\cdot d\mathbf{S} = 4\pi C_1 = \kappa M$, giving $C_1 = GM/c^2$. $\square$
\end{proof}

The $1/r$ decay is a geometric consequence of $d=3$: the Laplace equation Green's function in $d$ dimensions decays as $r^{-(d-2)}$. Gravity's inverse-square law is a statement about the dimensionality of space, not a separate postulate.

\subsection{Scale Comparison}

\begin{table}[h]
\caption{$q$ at surface for systems spanning the mass range.}
\begin{tabular}{lccc}
\toprule
System & $M$ [kg] & $R$ [m] & $q(R)$ \\
\midrule
Proton & $1.7\times 10^{-27}$ & $8.4\times 10^{-16}$ & $\approx 1-10^{-39}$ \\
Human & $70$ & $0.5$ & $\approx 1-10^{-25}$ \\
Earth & $6.0\times 10^{24}$ & $6.4\times 10^6$ & $0.9999999993$ \\
Sun & $2.0\times 10^{30}$ & $7.0\times 10^8$ & $0.9999979$ \\
Neutron star & $1.4M_\odot$ & $10^4$ & $0.83$ \\
Black hole (at $r_s$) & $10M_\odot$ & $3\times 10^4$ & $2/3\to 0$ \\
\bottomrule
\end{tabular}
\end{table}

Surprisingly, Earth is an extraordinarily quantum object in the DGF sense---less than $10^{-7}$\% of its Planck cells are occupied. Yet this tiny deviation, integrated over Earth's volume and amplified by the $1/q^2$ diffusivity, produces $g=9.8$~m/s$^2$.

\subsection{What is Derived and What is Not}

We clearly distinguish derived structure from calibrated constants:

\textbf{Derived:} The Laplace/Poisson \emph{form} of the field equation, the $1/r$ decay, Gauss's law structure, and the structural identity with Newtonian gravity.

\textbf{Not derived:} The numerical value of $G$, the mass-information conversion factor $\eta$, the propagation speed (currently parabolic, not hyperbolic), spacetime curvature, and gravitational waves.

The framework does not claim to replace general relativity. It claims that the \emph{static Newtonian limit} of gravity is a consequence of information dynamics under finite capacity---a result that no other framework derives from axioms as minimal as A1--A3.

% ============================================================
\section{Physical Interpretation}
\label{sec:interpretation}

\subsection{$q$ as a Coherence Measure}

The $q$ parameter admits a precise mapping to standard quantum mechanics. In the pointer basis $\{|i\rangle\}$ selected by the system-environment interaction Hamiltonian, the $\ell_1$-norm of coherence~\cite{Baumgratz2014}:
\begin{equation}
C_{\ell_1}(\rho) = \sum_{i\neq j}|\rho_{ij}|
\end{equation}
quantifies the total magnitude of off-diagonal elements. The normalized version:
\begin{equation}
q(\rho) = \frac{C_{\ell_1}(\rho)}{\|\rho\|_1} = \frac{\sum_{i\neq j}|\rho_{ij}|}{\sum_{i,j}|\rho_{ij}|}
\end{equation}
maps $q=1$ (all elements have equal magnitude, maximum coherence) to the DGF ``fully quantum'' limit, and $q=0$ (diagonal density matrix, zero coherence) to the DGF ``fully classical'' limit.

This mapping satisfies the required consistency checks:
\begin{enumerate}
\item \textbf{Monotonicity:} Under phase damping $\rho_{ij}\to (1-p)\rho_{ij}$ for $i\neq j$, $q$ strictly decreases for $p>0$.
\item \textbf{Absorbing boundary:} $q=0$ (diagonal state) is invariant under any decoherence channel.
\item \textbf{Complementarity with S1:} The reflux bound constrains information \emph{capacity}, while standard decoherence constrains information \emph{rates}. The bound is not derivable from standard Lindblad theory.
\end{enumerate}

\subsection{The Information Flow Story}

The framework's central physical narrative can be stated in one paragraph:

A pure quantum universe begins with $q\approx 1$ (all cells vacant). Initial inhomogeneities create $q$-gradients. Information flows from low-$q$ (occupied) to high-$q$ (vacant) regions driven by the Shannon entropy gradient. Regions that accumulate enough occupation approach $q\approx 0$---a self-organized critical state where any perturbation triggers cascaded overflow. These critical regions are what we call ``mass.'' They do not merely sit there---their criticality makes them powerful sources of $q$-field gradients, extending to infinity as $1/r$. Surrounding regions, maintained at $q\approx 1/2$ by the entropic preference, receive this information flow while retaining enough vacant capacity to sustain quantum coherence. The classical world is not the quantum world's corpse---it is the quantum world's engine.

\subsection{Comparison with Standard Approaches}

\textbf{Zurek's quantum Darwinism:} Standard decoherence requires an environment to select pointer states and destroy coherence. DGF shows that the same phenomena emerge from capacity asymmetry alone---even a perfectly isolated system develops a preferred basis through its own interaction history. The difference is experimentally testable: in DGF, the decoherence rate depends on the \emph{capacity ratio} $q_S/q_E$, not just on the environmental spectral density.

\textbf{Penrose/Di\'osi gravitational decoherence:} These models attribute decoherence to gravitational self-energy, predicting $\tau\propto 1/\rho^{1/3}$ (denser objects decohere \emph{faster}). DGF with Process~B (spontaneous occupation) predicts the opposite: $\tau\propto\rho$ (denser objects decohere \emph{slower}, because equal mass means fewer Planck cells). We note that Process~B is a conjecture, not a theorem, and we do not claim the $\tau\propto\rho$ prediction as a firm result of the framework.

% ============================================================
\section{Experimental Signature}
% ============================================================

\subsection{Information Reflux in Circuit QED}

The S1 reflux bound provides the cleanest experimental discriminator between DGF and standard open quantum systems theory.

\textbf{System:} A transmon qubit (system~S, $N_S\approx 1$ effective bit at low temperature) coupled to a readout cavity or an array of $N_E$ ancilla qubits (environment~E). At dilution refrigerator temperatures ($T_{\rm eff}\sim 35$~mK), thermal excitations are negligible.

\textbf{Standard Lindblad prediction:} At zero temperature, the Lindblad master equation with cavity damping gives $P(|e\rangle\to|g\rangle\to|e\rangle)=0$---once the qubit decays, it never re-excites.

\textbf{DGF prediction:} The finite information capacity of the environment imposes an upper bound on how completely it can absorb the qubit's information. The reflux probability is:
\begin{equation}
P_{\rm reflux} \le \frac{q_S}{q_E} \approx \frac{1}{N_E+1},
\end{equation}
where the approximation holds for $q_S\approx 1$ (qubit initially excited) and $q_E\approx 1$ (environment initially in vacuum).

\textbf{Proposed experimental protocol:}
\begin{enumerate}
\item Prepare transmon in $|e\rangle$, environment in vacuum.
\item Let the system evolve for time $T$ (spanning $T_1$ to several $T_1$).
\item Perform quantum state tomography to extract the trace distance $D(\rho(t),\rho(0))$.
\item Identify reflux events as increases in the trace distance after an initial decrease.
\item Vary the effective environment size $N_E$ (by changing the number of coupled ancilla qubits or cavity modes).
\item Verify $P_{\rm reflux}\le 1/(N_E+1)$.
\end{enumerate}

\textbf{Numerical estimates:} Using published transmon parameters (Yale 3D cavity: $T_1=80~\mu$s, $g/2\pi=160$~MHz, $T_{\rm eff}=35$~mK, residual excitation $0.1\%$)~\cite{Jin2015}. For $N_E=30$--$200$ (achievable with ancilla qubit arrays), DGF predicts $P_{\rm reflux}=0.5$--$3.2\%$, distinguishable from the quasiparticle background ($\sim 0.1\%$) with $10^4$--$10^6$ repetitions (requiring $0.9$--$170$~seconds of integration time).

\textbf{Falsification criterion:}
\begin{itemize}
\item If $P_{\rm reflux}=0$ at $T\ll\hbar/k_BT$: DGF is falsified; standard Lindblad is supported.
\item If $P_{\rm reflux}>0$ but $\le q_S/q_E$: Standard Lindblad is challenged; DGF survives.
\item If $P_{\rm reflux}>q_S/q_E$: DGF is falsified.
\end{itemize}

\textbf{Honest caveats:} (1)~The mapping between $q$ and experimental observables (open problem~O5) uses the operational proxy $q_S\approx T_2/(T_1+T_2)$, which is not a strict derivation. (2)~In the closed-system regime, DGF and standard quantum mechanics predict the same microcanonical typicality value $P(|e\rangle)=1/(N_E+1)$; the discriminating power comes from the \emph{upper bound} structure unique to DGF. (3)~The experimental protocol requires excellent isolation from thermal and quasiparticle backgrounds. Current technology (IBM Q, Yale 3D cavities) is approaching but has not yet demonstrated the required noise floor.

% ============================================================
\section{Discussion}
% ============================================================

\subsection{What the Framework Explains}

The three axioms A1--A3, together with the discrete-to-continuum limit, explain:
\begin{enumerate}
\item \textbf{Why the quantum-classical separation is directional.} $q\approx 0$ is a critical state (self-organized criticality); $q\approx 1/2$ is the entropic attractor. The separation is the coexistence of these two regimes under the $q$-field dynamics---not a mystery, but a phase coexistence.
\item \textbf{Why irreversibility depends on environment size.} Theorem~3 separates the relaxation timescale (controlled by $N_E$) from the recurrence timescale (controlled by $N_S$). ``Irreversibility'' is fast relaxation relative to observation time, not fundamental loss of information.
\item \textbf{Why the gravitational field equation has the form it does.} The static limit of the $q$-field is the Laplace equation $\nabla^2(1/q)=0$, identical in structure to $\nabla^2\Phi=0$. The $1/r$ decay follows from $d=3$. The equation structure is derived; only the coupling constant $G$ is calibrated.
\item \textbf{Why the classical world does not subsume the quantum world.} The entropic preference for $q\approx 1/2$ maintains a reservoir of quantum coherence. The classical cores ($q\approx 0$) and the quantum background ($q\approx 1/2$) are two phases of the same $q$-field, locked in dynamic equilibrium by information flow.
\end{enumerate}

\subsection{What the Framework Does Not Explain}

\begin{enumerate}
\item \textbf{The numerical values of $c$, $\hbar$, and $G$.} These are Level~4 (contingent) parameters. The framework explains \emph{structures} (Laplace equation, $1/r$ decay, phase coexistence), not \emph{specific values}.
\item \textbf{Why $d=3$.} Spatial dimensionality is an input. The $1/r$ decay of gravity would be $1/r^{d-2}$ in $d$ dimensions---this is a testable link between gravity and dimension, not a derivation of dimension.
\item \textbf{Quantum field theory and the standard model.} The cell model is a discrete classical information system. Bridging to relativistic QFT requires additional structure not contained in the current framework.
\item \textbf{The exact value of the cosmological constant.} The $10^{-5}$ CMB residual and the $\Lambda\sim 10^{-120}$ hierarchy are not addressed.
\item \textbf{Process~B and spontaneous occupation.} Whether an additional mechanism beyond $q$-field diffusion is needed to drive $q\to 0$ remains an open question. The $\tau\propto\rho$ prediction depends on Process~B and should be treated as conditional.
\end{enumerate}

\subsection{Relation to Bassani \& Magueijo}

Bassani and Magueijo~\cite{Bassani2025} proposed a natural selection mechanism for physical laws, using unimodular gravity and absorbing Markov chains. Their framework answers \emph{how} a universe with fixed laws can emerge from random mutations. Ours answers a complementary question: \emph{why}, given minimal information-theoretic axioms, the separation of scales and the structure of gravity take the form they do. The two frameworks are not competitors---they address different layers of the same problem. A synthesis, in which DGF provides the information-theoretic ``environment'' for BM-style natural selection, would be a natural next step.

% ============================================================
\section{Conclusion}
% ============================================================

This paper has shown that three axioms about information---distinguishability, bounded capacity, and locality---are sufficient to explain the structural features of the quantum-classical separation and the form of the gravitational field equation. We have not derived the universe. We have shown that certain features of the universe that appear contingent are, in fact, necessary consequences of what it means for information to exist and flow under finite capacity.

The framework makes one clean experimental prediction: an information reflux bound $P_{\rm reflux}\le q_S/q_E$, testable in circuit QED at zero temperature. Verification of this bound would establish information capacity as a genuine physical constraint on quantum dynamics, distinct from standard decoherence theory.

The deepest insight of this work is perhaps the simplest: the classical world is not the quantum world's graveyard. It is the quantum world's engine---a self-organized critical state whose static geometry is what we measure as gravity, and whose information outflow sustains the quantum coherence of everything around it.

\begin{thebibliography}{99}

\bibitem{Zurek2003} W.H.~Zurek, Rev.~Mod.~Phys.~\textbf{75}, 715 (2003).

\bibitem{Schlosshauer2007} M.~Schlosshauer, \emph{Decoherence and the Quantum-to-Classical Transition} (Springer, 2007).

\bibitem{Zurek2009} W.H.~Zurek, Nat.~Phys.~\textbf{5}, 181 (2009).

\bibitem{Bak1987} P.~Bak, C.~Tang, and K.~Wiesenfeld, Phys.~Rev.~Lett.~\textbf{59}, 381 (1987).

\bibitem{Baumgratz2014} T.~Baumgratz, M.~Cramer, and M.B.~Plenio, Phys.~Rev.~Lett.~\textbf{113}, 140401 (2014).

\bibitem{Jin2015} X.Y.~Jin \emph{et al.}, Phys.~Rev.~Lett.~\textbf{114}, 240501 (2015).

\bibitem{Bassani2025} P.M.~Bassani and J.~Magueijo, Phys.~Rev.~D~\textbf{111}, 103529 (2025).

\bibitem{Cover2006} T.M.~Cover and J.A.~Thomas, \emph{Elements of Information Theory} (Wiley, 2006).

\bibitem{Chiribella2011} G.~Chiribella, G.M.~D'Ariano, and P.~Perinotti, Phys.~Rev.~A~\textbf{84}, 012311 (2011).

\bibitem{Malament1977} D.~Malament, J.~Math.~Phys.~\textbf{18}, 1399 (1977).

\bibitem{Penrose1996} R.~Penrose, Gen.~Rel.~Grav.~\textbf{28}, 581 (1996).

\bibitem{Diosi1989} L.~Di\'osi, Phys.~Rev.~A~\textbf{40}, 1165 (1989).

\bibitem{Vopson2023} M.~Vopson, AIP Advances \textbf{13}, 10520 (2023).

\bibitem{Verlinde2011} E.~Verlinde, JHEP \textbf{04}, 029 (2011).

\bibitem{Sorkin1991} R.D.~Sorkin, in \emph{Relativity and Gravitation} (1991).

\bibitem{Wheeler1990} J.A.~Wheeler, in \emph{Complexity, Entropy, and the Physics of Information} (1990).

\end{thebibliography}

\end{document}
